\documentclass[spanish,letterpaper,oneside]{article}

\def\documenttitle {Cálculo Integral}
\def\documentsubtitle {Portafolio de evidencias}
\def\documentsubject {Ingeniero Agrónomo}
\def\documentauthor {Diego de la Cruz}
\def\coursename {Cálculo Integral}
\def\coursecode {Clave por registrar}
\def\universityname {Universidad Autónoma de Nuevo León}
\def\universityfaculty {Facultad de Agronomía}
\def\universitydepartment {Ingeniero Agrónomo}
\def\universitydepartmentimage {UANL/assets-uanl/logo-uanl.png}
\def\universitydepartmentimagecfg {width=3.2cm}
\def\universitylocation {General Escobedo, Nuevo León, México}
\def\authortable {
  \begin{tabular}{ll}
    \textbf{Alumno:} & Diego de la Cruz \\
    Matrícula: & Por registrar \\
    Grupo: & Por registrar \\
    Programa: & Ingeniero Agrónomo \\
    Unidad de aprendizaje: & \coursename \\
    Docente: & Por registrar \\
    & \\
    \multicolumn{2}{l}{Fecha de realización: \today} \\
    \multicolumn{2}{l}{\universitylocation}
  \end{tabular}
}

\input{template}

\begin{document}
\templatePortrait
\templatePagecfg
\onehalfspacing

\begin{abstractd}
Este portafolio reúne las evidencias de Cálculo Integral de Diego de la Cruz en el programa de Ingeniero Agrónomo de la UANL. Cada reporte presenta una síntesis de fórmulas y procedimientos y la resolución completa, justificada y ordenada de los ejercicios asignados.
\end{abstractd}

\templateIndex
\templateFinalcfg

\section{Evidencias}
\begin{enumerate}
  \item Actividad ponderada: integración por partes (6 puntos).
  \item Evidencia 1: aprendizaje basado en problemas (20 puntos).
  \item Evidencia 2: integración de funciones racionales (15 puntos).
  \item Actividad ponderada 2.2: modelos de ganancia de peso (6 puntos).
  \item PIA de Cálculo Integral: integrales algebraicas y racionales (30 puntos).
\end{enumerate}

\nocite{*}
\bibliography{calculo-integral}
\end{document}
